% Fully read excerpts from https://raw.githubusercontent.com/openai/math/fd4aeeb2ee4fc729c18d98444fed42fd0529eeeb/preprints/The-complete-Crouzeix-theorem-September-23-2026/build/main.tex

\subsection*{Coordinates and evaluation}
Write $\D=\{z\in\C:|z|<1\}$ and $\T=\partial\D$, with normalized
angular measure $d\sigma(e^{i\theta})=d\theta/(2\pi)$. For
$A\in M_n(\C)$ its numerical range is
\[
 W(A)=\{x^*Ax:x\in\C^n,\ x^*x=1\}.
\]
All operator norms are Hilbert-space operator norms, and tensor products are
Hilbert tensor products with the base space first. For a matrix-valued function $v$
holomorphic near the spectrum, $v[A]$ denotes holomorphic evaluation:
on a Jordan block $J=\beta I+N$ of size $s$ it is
$\sum_{j=0}^{s-1}N^j\otimes v^{(j)}(\beta)/j!$, followed by the
change of basis on the first factor. In particular,
$v[A]=\sum_jA^j\otimes B_j$ when $v(z)=\sum_jB_jz^j$.
Inner products are conjugate-linear in the first argument, and $U\preceq V$
means that $V-U$ is positive semidefinite. The symbol ${}^*$ denotes the
adjoint on a base or coefficient Hilbert space, ${}^{\mathsf T}$ ordinary
transpose, and ${}^\dagger$ the adjoint of an operator on a function space.

Call $\Omega$ \emph{admissible} when it is a bounded open convex subset of
$\C$ with regular real-analytic Jordan boundary.

\Needspace{9\baselineskip}
\begin{lemma}[Conformal coordinates]\label{lem:conformal-coordinates}
For every admissible $\Omega$ there is a conformal bijection
$f:\Omega\to\D$ such that $f$ and $f^{-1}$ extend univalently
through $\overline\Omega$ and $\overline\D$, respectively. There is also
a conformal bijection $G:\{|t|>1\}\to\C\setminus\overline\Omega$
that extends univalently to
$\{|t|>r\}\cup\{\infty\}$ for some $0<r<1$, as a sphere map, with a
simple pole at infinity and $G(\T)=\partial\Omega$. On $\T$, $tG'(t)$ is an outward normal and
$\psi=f\circ G$ is an orientation-preserving smooth circle diffeomorphism.
\end{lemma}
The complete local Riemann-mapping, harmonic-reflection, and collar argument
is in Appendix~\ref{sec:conformal-appendix}.


Fix such an $\Omega$ with $W(A)\subset\Omega$, and choose its coordinates
$f,G$ from Lemma~\ref{lem:conformal-coordinates}. Put $T=f(A)$ and define
\begin{equation}\label{eq:minimum}
 \kappa^2=\min\{\tau\in\R:I\preceq H\preceq\tau I,
                  \ T^*HT\preceq H\text{ for some }H=H^*\}.
\end{equation}
A strictly feasible pair means that $H-I$, $\tau I-H$, and
$H-T^*HT$ are all positive definite. Existence and attainment of the
minimum are part of the first conclusion.

\begin{theorem}[The optimal metric]\label{thm:structural}
For every admissible $\Omega$, every finite $n\ge1$, every
$A\in M_n(\C)$ with $W(A)\subset\Omega$, and the interior conformal
coordinate $f$ above, the problem \eqref{eq:minimum} is strictly feasible
and has an attained minimum $1\le\kappa^2\le4$. If $H$ is a minimizing
metric, then
\[
 S=H^{1/2},\qquad A'=SAS^{-1},\qquad D=Sf(A)S^{-1}=f(A')
\]
satisfy
\[
 D^*D\preceq I,\qquad \rho(D)<1,
 \qquad \|S\|\|S^{-1}\|=\kappa\le2.
\]
\end{theorem}

\begin{theorem}[A common positive representation]\label{thm:representation}
With the same $\Omega,A,f,G$ and any chosen minimizing metric in
Theorem~\ref{thm:structural}, there is one continuous function
$\Lambda:\T\to M_n(\C)$ satisfying
\[
 \Lambda(t)\succeq0,\qquad \int_\T\Lambda(t)\,d\sigma(t)=I_n,
\]
such that, for every positive integer $m$ and every $M_m(\C)$-valued
function $v$ holomorphic on a neighborhood of $\overline\Omega$,
\begin{equation}\label{eq:representation-target}
 v[A']=\int_\T\Lambda(t)\otimes v(G(t))\,d\sigma(t),
 \qquad
 \|v[A]\|\le\kappa\max_{\overline\Omega}\|v\|.
\end{equation}
The same $S$, $\Lambda$, and $\kappa$ work for all $m$ and $v$.
\end{theorem}



% --- next excerpt ---

\subsection{A metric for the disk coordinate}
The evaluation rules used below follow directly from the finite Jordan
Taylor formula. Its product rule preserves the order of matrix coefficients;
conjugating the Jordan form proves similarity on the base factor. For scalar
composition, substitute the nilpotent matrix $f(J)-f(\beta)I$ into the
Taylor series of the outer function at $f(\beta)$. The finite Taylor
composition identity gives $(v\circ f)[J]=v[f(J)]$, and the same formula
gives $\spec(f(A))=f(\spec(A))$. No diagonalizability is assumed.
The Jordan resolvent expansion and Cauchy's derivative formula give the
equivalent contour-integral evaluation.

Since $\spec(A)\subset W(A)\subset\Omega$, spectral mapping gives
$\rho(T)<1$. The positive-metric feasibility criterion is classical
\cite[Theorem 1]{Stein1952}; we construct the strictly feasible witness.
Choose $\rho(T)<\eta<1$. Jordan normal form supplies
$C_\eta$ with $\|T^j\|\le C_\eta\eta^j$. Hence
\[
 H_0=\sum_{j\ge0}T^{*j}T^j
\]
converges in norm, $H_0\succeq I$, and $H_0-T^*H_0T=I$.
For $c>1$ and $\tau>\|cH_0\|$, the pair $(\tau,cH_0)$ has all
three matrix slacks in \eqref{eq:minimum} positive definite. This is
strict feasibility of the optimization problem. Feasible pairs below any
fixed feasible upper bound form a closed bounded subset of a finite-dimensional
space, so a minimum is attained. Every feasible $\tau$ is at least one.

Fix a minimizing $H$ and put $S=H^{1/2}$. Congruence gives
$D^*D=S^{-1}T^*HTS^{-1}\preceq I$; similarity gives $\rho(D)<1$.
The inequalities $I\preceq H\preceq\kappa^2I$ give
$\|S\|\|S^{-1}\|\le\kappa$. The optimizer itself may have zero
slacks. In particular $\|D\|=1$ is allowed.

The minimum also measures the best condition number of a similarity to
a contraction. If $R$ is invertible and $\|RTR^{-1}\|\le1$, then
$H_R=R^*R$ satisfies $T^*H_RT\preceq H_R$. Dividing $H_R$ by its
least eigenvalue gives a feasible metric with
$\tau=\|R\|^2\|R^{-1}\|^2$. Consequently
$\kappa\le\|R\|\|R^{-1}\|$. Applying this to the minimizing
$S=H^{1/2}$, for which the reverse inequality was just proved, gives
\begin{equation}\label{eq:optimal-similarity}
 \kappa=\min\{\|R\|\|R^{-1}\|:
 R\text{ invertible},\ \|RTR^{-1}\|\le1\}.
\end{equation}
The minimum is attained by the positive matrix $S$, with
$\|S\|\|S^{-1}\|=\kappa$. This identity holds for any fixed finite
$T$ with an attained feasible metric problem; the stability assumption
above supplies such feasibility. The bound $\kappa\le2$ proved below
uses the additional conformal and numerical-range hypotheses.

\subsection{One density for every coefficient size}
\begin{proposition}[One density for every coefficient size]\label{prop:density}
For the fixed minimizing metric, the representation
\eqref{eq:representation-target} holds, before any bound on $\kappa$ is
proved, with
\begin{align}
 \Phi_D(\zeta)&=(I-D/\zeta)^{-1}+[(I-D/\zeta)^{-1}]^*-I,
                                                        \label{eq:phi}\\
 j(t)&=\frac{tG'(t)f'(G(t))}{\psi(t)}>0,
 &\Lambda(t)&=j(t)\Phi_D(\psi(t)).                      \label{eq:lambda}
\end{align}
\end{proposition}
The disk-resolvent formula is the positive representation of
Delyon and Delyon \cite[Section 2, equation (7)]{DelyonDelyon1999};
the common compression also belongs to the dilation tradition
\cite[Theorem I]{SzNagy1953}. We express the density in the exterior
parameter and prove every coefficient-size claim.
\begin{proof}
For $|\zeta|=1$, $E_\zeta=I-D/\zeta$ is invertible and
\[
 E_\zeta^*\Phi_D(\zeta)E_\zeta=I-D^*D\succeq0.
\]
Thus $\Phi_D$ is positive and continuous. The strict spectral-radius
bound, rather than a strict norm bound, implies $\sum_{j\ge0}\|D^j\|<\infty$.
Consequently the following series is absolutely uniformly convergent on
the circle:
\[
 \Phi_D(\zeta)=I+\sum_{j\ge1}(\zeta^{-j}D^j+\zeta^jD^{*j}),
 \qquad \int_\T\zeta^j\Phi_D(\zeta)\,d\sigma(\zeta)=D^j
 \quad(j\ge0).
\]
For any $m$ and any $M_m$-valued $h$ holomorphic near $\overline\D$,
integration of its uniformly convergent Taylor series gives
\[
 h[D]=\int_\T\Phi_D(\zeta)\otimes h(\zeta)\,d\sigma(\zeta).
\]
Apply this with $h=v\circ f^{-1}$. The conformal collar makes $h$
holomorphic near $\overline\D$, and scalar composition gives $h[D]=v[A']$.
Locally writing $\psi(e^{i\theta})=e^{i\eta(\theta)}$ shows that
$j(e^{i\theta})=\eta'(\theta)>0$. Changing variables
$\zeta=\psi(t)$ proves the first identity in
\eqref{eq:representation-target}. The zeroth moment gives
$\int\Lambda\,d\sigma=I$.

For each $m$, define
\[
 V_m:\C^n\otimes\C^m\longrightarrow
 L^2(\T,d\sigma;\C^n\otimes\C^m),\qquad
 (V_m b)(t)=(\Lambda(t)^{1/2}\otimes I_m)b.
\]
Normalization makes $V_m$ an isometry. If $M_v$ multiplies by
$I_n\otimes v(G(t))$, then
\[
 v[A']=V_m^\dagger M_vV_m,\qquad
 \|v[A']\|\le\max_{\overline\Omega}\|v\|.
\]
Similarity on the first factor gives the remaining estimate with constant
$\kappa$. Neither the definition of $S$ nor that of $\Lambda$ depends
on $m$ or $v$.
\end{proof}

\subsection{Endpoint separation and an exact extremal}
To bound the metric, we need equality for one auxiliary matrix function.
The construction passes from real separation to a dual balance, then to
endpoint supports, and finally to a finite unitary realization.
Use the vectorization convention
\[
 \vecop(X)=\sum_{j=1}^nXe_j\otimes e_j,\qquad
 \|\vecop(X)\|^2=\tr(XX^*).
\]

\begin{proposition}[Endpoint extremal]\label{prop:extremal}
If $\kappa>1$, there are $X,Y\in M_n(\C)$ with
\[
 \tr(XX^*)=\tr(YY^*)=1,\qquad SX=\kappa X,\qquad SY=Y,
 \qquad X^*Y=0,
\]
and an $M_n$-valued rational function $F_\D$, holomorphic on a neighborhood
of $\overline\D$ and contractive on $\overline\D$, such that, for
$x=\vecop(X)$ and $y=\vecop(Y)$,
\[
 F_\D[D]x=y.
\]
In particular $F_\Omega=F_\D\circ f$ is holomorphic near
$\overline\Omega$ and satisfies $F_\Omega[A']x=y$. Its boundary trace
$F=F_\Omega\circ G$ is contractive.
\end{proposition}
\begin{proof}
We first obtain positive matrices $X_0$ and $Y_0$ supported on the
$\kappa^2$ and $1$ eigenspaces of $H$, respectively. A third positive
matrix $Z_0$ will express their difference as $Z_0-TZ_0T^*$; this balance
will provide the equal Gram matrices needed for the unitary realization.

\emph{Separation, including a possible contact point.}
In the real vector space $\R\oplus\operatorname{Herm}_n^3$, the affine set
\[
 \mathcal A_T=\{(s,J-I,sI-J,J-T^*JT):s\in\R,\ J=J^*\}
\]
is disjoint from the open convex set
\[
 \mathcal O=\{(s,U,V,W):s<\kappa^2,\ U,V,W\succ0\}.
\]
Here is the finite-dimensional separation argument used for these sets;
the general convex-separation framework is treated in
\cite[Section 11, Theorem 11.2]{Rockafellar1970}.
Write $\mathcal A_T=a_0+L$, and let $\pi$ project orthogonally onto
$L^\perp$. The set $\mathcal K=\pi(\mathcal O-a_0)$ is nonempty, open
in $L^\perp$, and convex, and it does not contain zero. If
$0\notin\overline{\mathcal K}$, a nearest point $z_0$ in the closed convex
set $\overline{\mathcal K}$ satisfies
$\langle z_0,c\rangle\ge\|z_0\|^2>0$ for every $c\in\mathcal K$.

If $0\in\overline{\mathcal K}$, fix $c_*\in\mathcal K$ and set
$z_j=-c_*/j$. Each $z_j$ is outside $\overline{\mathcal K}$: otherwise
zero would lie on the open segment joining an interior point $c_*$ to a
point of the closure, and hence would belong to $\mathcal K$. To verify
that segment fact, take a ball about $c_*$ inside $\mathcal K$, approximate
the other endpoint by points of $\mathcal K$, and use convexity; a ball
about each fixed interior point of the limiting segment is still contained
in $\mathcal K$. Let $y_j$ be nearest to $z_j$ in
$\overline{\mathcal K}$. Since zero belongs to this closed set,
$\|y_j-z_j\|\le\|z_j\|$, so $y_j\to0$. A subsequence of
$\nu_j=(y_j-z_j)/\|y_j-z_j\|$ tends to a unit vector $\nu$.
The nearest-point inequality gives
$\langle\nu_j,c-y_j\rangle\ge0$, hence
$\langle\nu,c\rangle\ge0$ for every $c\in\mathcal K$.
In the first case use $\nu=z_0$. In either case
$\ell(v)=-\langle\nu,\pi v\rangle$ is nonzero and
$\ell(a)\ge\ell(o)$ for all $a\in\mathcal A_T$, $o\in\mathcal O$.

Because $\mathcal A_T$ is affine, $\ell$ is constant on it. Boundedness
above on $\mathcal O$ forces the signs
\[
 \ell(s,U,V,W)=\lambda s-\tr(Y_0U)-\tr(X_0V)-\tr(Z_0W),
 \quad\lambda\ge0,\quad X_0,Y_0,Z_0\succeq0.
\]
If $\lambda=0$, strict feasibility produces a negative value on
$\mathcal A_T$, while separation requires a value at least
$\sup_{\mathcal O}\ell=0$. Thus $\lambda>0$; scale it to one.
Constancy on $\mathcal A_T$ then gives
\begin{equation}\label{eq:dual}
 X_0-Y_0=Z_0-TZ_0T^*,\qquad \tr X_0=1.
\end{equation}
The constant value is $\tr Y_0\ge\kappa^2$. At $(\kappa^2,H)$ it is
\[
 \kappa^2-\tr(Y_0(H-I))-\tr(X_0(\kappa^2I-H))
            -\tr(Z_0(H-T^*HT))\le\kappa^2.
\]
Each subtracted trace is nonnegative, so all vanish and
$\tr Y_0=\kappa^2$.

\emph{Endpoint support and balance.}
Write
$\|B\|_{\HS}=(\tr(B^*B))^{1/2}$ for the unnormalized
Hilbert--Schmidt norm. For positive matrices $B,C$,
$\tr(BC)=\|C^{1/2}B^{1/2}\|_{\HS}^2$; vanishing implies that the
range of $B$ lies in the kernel of $C$. Therefore $X_0$ is supported on
the $\kappa^2$ eigenspace of $H$, and $Y_0$ on its $1$ eigenspace.

Choose square factors and a positive matrix $Q$ by
\[
 XX^*=\kappa^{-2}SX_0S,\quad
 YY^*=\kappa^{-2}SY_0S,\quad Q=\kappa^{-2}SZ_0S.
\]
The support and trace relations give the asserted norms and endpoint
identities. Their column ranges are orthogonal because $\kappa>1$,
so $X^*Y=0$ as a matrix. Congruence of \eqref{eq:dual} gives
\begin{equation}\label{eq:balance}
 XX^*-YY^*=Q-DQD^*.
\end{equation}

\emph{Unitary realization and transpose.}
The balance-to-isometry step is the classical interpolation mechanism of
\cite[Section~3 of the versioned preprint]{KatsnelsonKheifetsYuditskii1997}; its particular endpoint
data and finite-dimensional realization are established here.
Put $L=Q^{1/2}$. The rows $R_1=[L\ Y]$ and $R_2=[DL\ X]$ have
equal products with their adjoints. The map
$R_2^*v\mapsto R_1^*v$ is therefore a well-defined isometry of their
ranges. Extend it to a unitary using equal-dimensional orthogonal
complements, then take its adjoint. This yields
\[
 [L\ Y]=[DL\ X]
 \begin{pmatrix}a&b\\c&d\end{pmatrix},\qquad
 U=\begin{pmatrix}a&b\\c&d\end{pmatrix}\text{ unitary}.
\]
Thus $L=DLa+Xc$ and $Y=DLb+Xd$. Since $\|a\|\le1$ and
$\rho(D)<1$, iteration gives the norm-convergent identity
\begin{equation}\label{eq:realization-series}
 Y=Xd+\sum_{j\ge0}D^{j+1}Xca^jb.
\end{equation}
The transfer-function norm identity is the unitary-colligation identity
of \cite[Section 3.1, equations (3.1)--(3.4)]{Brodskii1978}.
Define, for $|\zeta|<1$,
\[
 R(\zeta)=d+\zeta c(I-\zeta a)^{-1}b,
 \qquad F_\D(\zeta)=R(\zeta)^{\mathsf T}.
\]
For $e\in\C^n$ and $v=(I-\zeta a)^{-1}be$, unitarity gives
\[
 U\binom{\zeta v}{e}=\binom{v}{R(\zeta)e},\qquad
 \|e\|^2-\|R(\zeta)e\|^2=(1-|\zeta|^2)\|v\|^2\ge0.
\]
Hence $R$ is contractive. Ordinary transpose preserves its norm and
holomorphy. Its rational entries are bounded in $\D$, so every apparent
pole on $\T$ is removable: a genuine pole would diverge on an interior
radial approach. After cancellation, finitely many remaining poles stay
away from $\overline\D$. This supplies a common neighborhood of that
closed disk on which $F_\D$ is holomorphic; continuity retains
contractivity on the circle.

Finally,
\[
 (D^k\otimes B)\vecop(X)=\vecop(D^kXB^{\mathsf T}).
\]
The Taylor coefficients of $F_\D$ are the transposes of $d$ and
$ca^jb$. Substitution in \eqref{eq:realization-series} proves
$F_\D[D]x=y$. The conformal collar and scalar composition give the
claims for $F_\Omega$. Rationality is asserted in the disk variable;
$F_\Omega$ is only required to be holomorphic near $\overline\Omega$.
\end{proof}

